% Shared exercise chunk: l2-unknown-sampling-law
\begin{exercisechunk}
\item \textbf{Two meanings of ``the learner does not know $Q$.''}
\label[exercise]{ex:unknown-sampling-law}
Assume $N\ge 2$ and $m\ge 1$, and compare
\begin{equation}
\min_{\cA}\max_{Q\in\cQ}\max_f R_Q(\cA,f)
\label{eq:absolute-distribution-free}
\end{equation}
with the requirement
\begin{equation}
\cA\in
\bigcap_{Q\in\cQ}
\operatorname*{arg\,min}_{\cB}\max_f R_Q(\cB,f).
\label{eq:simultaneous-optimality}
\end{equation}
\begin{enumerate}[(a)]
\item Show that the value of \cref{eq:absolute-distribution-free} is $1/2$.
\item Show that a learner that always predicts with a fair coin is optimal for
\cref{eq:absolute-distribution-free}, even though it ignores revealed labels.
\item Show that \cref{eq:simultaneous-optimality} forces the learner to return
the revealed label at every seen test index and to make an unbiased prediction
at every unseen test index, for every consistent labeled sample.
\item Does the learner in part~(b) satisfy \cref{eq:simultaneous-optimality}?
Give a sampling law witnessing your answer.
\end{enumerate}
\end{exercisechunk}
